
\section{Related Work}
\subsection{Reasoning with Language Models}

Chain-of-thought (CoT) prompting~\cite{wei2022chain_cot} demonstrates that eliciting intermediate reasoning steps significantly enhances large language models' performance on complex reasoning tasks. Building on this foundation, synthetic prompting~\cite{shao2023synthetic_synprompt} and STaR~\cite{zelikman2022star} explore automated generation of reasoning demonstrations: the former employs backward-forward construction to create high-quality exemplars, while the latter bootstraps rationale datasets from minimal seed examples and introduces rationalization, training models to justify correct answers post-hoc. Tree of Thoughts~\cite{yao2023tree_tot} extends CoT to breadth-first search over alternative reasoning branches, enabling backtracking and global solution evaluation. Quiet-STaR~\cite{zelikman2024quiet_quietstar} further generalizes this paradigm to learn reasoning from unstructured text across all token positions through parallel sampling.

To enhance reasoning reliability, recent work introduces verification mechanisms. Self-consistency~\cite{wang2022self_consistency} samples multiple reasoning paths and selects the most frequent answer, while self-verification~\cite{weng2023large_self_verification} and Self-Refine~\cite{madaan2023self_selfrefine} implement iterative refinement through self-critique. Addressing hallucinations in knowledge-intensive scenarios, CoT-RAG~\cite{li2025cot_rag} proposes knowledge graph-driven CoT generation with expert-provided decision trees that encapsulate domain reasoning logic, while self-refinement-enhanced knowledge retrieval~\cite{niu2024mitigating} identifies high-hallucination tokens via attribution analysis and corrects inaccuracies through targeted retrieval. DDCoT~\cite{zheng2023ddcot} extends this to multimodal reasoning by partitioning responsibilities between vision models and language models through duty-distinct prompting.
For self-improvement through data generation, impossible distillation~\cite{jung2023impossibledistill} filters low-quality samples generated from suboptimal models, while ReST~\cite{gulcehre2023reinforced_rest} employs offline reinforcement learning to iteratively improve policies from self-generated data. Despite these advances, efficiently obtaining high-quality trajectory and reasoning data for target domains remains challenging; we address this through a CSRS-centered data flywheel that ensures annotation quality while maintaining scalability. 

\subsection{GUI Agents}

Recent work, including UI-TARS~\cite{qin2025ui_uitars1}, UI-TARS-2~\cite{wang2025ui_uitars2}, OpenCUA~\cite{wang2025opencua}, GUI-Owl~\cite{ye2025mobile_mobileagentv3}, and UITron~\cite{zeng2025uitron}, systematically investigates data curation, cleaning pipelines, and training methodologies spanning pretraining and post-training phases. Concurrently, mainstream foundation models such as Claude series~\cite{claudeopus45}, Seed series~\cite{guo2025seed1}, and Qwen series~\cite{yang2025qwen3} begin natively integrating GUI interaction capabilities, underscoring the strategic importance of this research direction. Most GUI agents follow the ReAct paradigm~\cite{yao2022react}, observe, reason via CoT, then act, with recent systems extending this through multi-agent collaboration and hierarchical architectures. Notable examples include Mobile-Agent-v3~\cite{ye2025mobile_mobileagentv3}, which introduces specialized agents (Manager, Worker, Reflector, Notetaker) for knowledge evolution and reflective reasoning, and other systems~\cite{liu2025pc,song2025coact1,zhang2025ufo2,agashe2025agents2} that explore hierarchical planning, code-visual hybrid control, and compositional agent coordination. We propose the Step-GUI series models, achieving state-of-the-art performance among models of similar size, while Step-GUI-4B can run fully locally on consumer-grade hardware, enabling deployment in privacy-sensitive scenarios. 

To evaluate these advances, the community develops both static benchmarks including ScreenSpot~\cite{cheng2024seeclick}, ScreenSpot-v2~\cite{wu2024atlas}, ScreenSpot-Pro~\cite{li2025screenspotpro}, FuncPred~\cite{li2025autogui}, MoTIF~\cite{burns2022dataset_motif}, RefExp~\cite{rastogi2021uibert_refexp}, LlamaTouch~\cite{zhang2024llamatouch}, VWB-AG~\cite{liu2024visualwebbench_vwb}, and VWB-EG~\cite{liu2024visualwebbench_vwb}, that primarily assess element grounding and task planning accuracy through prediction-annotation matching, and interactive benchmarks such as AndroidWorld~\cite{rawles2024androidworld} and OSWorld~\cite{xie2024osworld} that enable agents to execute tasks within realistic or virtualized operating system environments with multi-dimensional task completion evaluation. In contrast to these efforts, we focus on scalable reasoning trajectory synthesis and efficient agent architecture design for real-world GUI automation scenarios.